\chapter[Challenges and Limitations]{Challenges and Limitations}
\label{cp:Challenges}

{
	\parindent0pt
	We acknowledge several challenges on the path to implementation.

	\begin{description}
		\item[Technical Challenges:]

		      \begin{itemize}
			      \item \textbf{Image Recognition Accuracy:} The model's accuracy is core to Ayang's integrity. Poor image quality could lead to errors. Our mitigation strategy is user education and aggressive model fine-tuning.
			  

			      \item \textbf{Malicious Use:} Users could attempt to cheat the system, for instance, by dumping their own trash to photograph. We plan countermeasures like GPS verification, time-stamping, and AI-powered anomaly detection to flag suspicious activity.

			      \item \textbf{Scalability and Cost:} User growth will increase API and storage costs. We will optimize the application and seek grants or partnerships to cover operational expenses.
		      \end{itemize}

		\item[Data and Ethical Challenges:]

		      \begin{itemize}
			      \item \textbf{Dataset Bias:} Publicly available trash datasets may not represent waste found in African communities. A model trained on Western data may be less accurate locally. This is why creating a custom, localized dataset is a core part of our methodology.

			      \item \textbf{Data Privacy:} The app collects location data, and we are committed to ethical handling. All user data will be anonymized for analysis, and our privacy policy will be transparent.
		      \end{itemize}

		\item[Logistical and Community Challenges:]

		      \begin{itemize}
			      \item \textbf{The "Last Mile" Problem:} The app assists with collection but not final disposal. If no bins are nearby, trash may just be moved to another pile. We will partner with local governments or waste companies to schedule pickups at designated collection points.

			      \item \textbf{Digital Divide:} Not everyone has a smartphone or affordable data, which may limit the app's reach. To promote inclusivity, we plan to explore offline functionality and community-based models where one person can document a group's work.
		      \end{itemize}
	\end{description}
}
